\documentclass[10pt,a4paper]{amsart}
\usepackage[margin=19mm]{geometry}
\usepackage{amsmath,amssymb,mathtools}
\usepackage[scr=rsfs,cal=euler]{mathalfa}
\usepackage{xfrac}
\usepackage[unicode,hidelinks,bookmarks=false]{hyperref}
\newcommand{\ie}{i.e.\ }
\newcommand{\cf}{cf.\ }
\newcommand{\resp}{respectively\ }
\newcommand{\Subsection}{\subsection}
\providecommand{\nobreakdash}{\nobreak\hbox{-}}
\setlength{\emergencystretch}{2em}
\multlinegap0pt
\usepackage{bm}
\usepackage{bbm}
\usepackage{dsfont}
\usepackage{xspace}
\usepackage{cancel}
\usepackage{mathtools}
\usepackage[shortlabels]{enumitem}
\usepackage{amsthm}
\makeatletter
\@ifundefined{lemma}{\newtheorem{lemma}{Lemma}}{}
\makeatother
\newcommand{\Ax}{\text{AX}}
\newcommand{\then}{\mathbin{\rightarrow}}
\newcommand{\tup}[1]{\langle #1 \rangle}
\renewcommand{\iff}{\mathbin{\leftrightarrow}}
\title[Stalnaker's logical problem of conditionals is unsolvable]{Stalnaker's logical problem of conditionals is unsolvable}
\author{Alexander W. Kocurek, James Walsh, Yale Weiss}
\date{Source: arXiv:2608.07387v1 (CC BY 4.0)}
\begin{document}
\maketitle

\noindent\emph{Source and licence.} Stalnaker's logical problem of conditionals is unsolvable. Version arXiv:2608.07387v1 (CC BY 4.0), distributed under the Creative Commons Attribution 4.0 International licence, as recorded in the arXiv licence metadata for this exact version and shown on the abstract page https://arxiv.org/abs/2608.07387v1. Full rights notice: licenses/arxiv-2608.07387-CC-BY-4.0.txt.\par\smallskip

\noindent\textit{Editorial scope.} Counted unit: the Lemma whose source label is \texttt{None} in the article (source0.tex lines 689--697, source label \texttt{None}), delivered together with the article's own complete original proof of it (source0.tex lines 699--780, one \texttt{proof} environment of 82 lines). No mathematics is written, repaired or re-derived: statement and proof are the article's own text line for line. Required context retained uncounted: ax:pred (lines 635--639); ax:zero (lines 635--639); lem:nat (lines 641--644); ax:add (lines 678--686); ax:cong (lines 678--686); ax:func (lines 678--686); ax:prod (lines 678--686); ax:unit (lines 678--686). Every definition, display and label of those lines is retained unchanged. Omitted from this package and not redistributed: 1--634, 640--640, 645--677, 687--688, 698--698, 781--925. Nothing inside a retained excerpt is omitted or altered: every retained body is delivered exactly as the article's lines are written, so no omission receipt of any kind accompanies this package. Citations: the retained excerpts carry no citation command at all, so no bibliography entry of the article is transcribed and no reference is redistributed. Reference adaptation: none was needed and none was made; every reference inside the delivered unit resolves inside it, so no unresolved reference or citation is produced. Printed slips kept literal: no printed word, symbol, hypothesis or number of the counted statement or of its proof was altered. Layout and class: the article's own class is \texttt{article} with packages including geometry, amsmath, amssymb, amsthm, enumitem, pifont, csquotes, tikz, caption, graphicx, biblatex; the wrapper is the lane's pinned \texttt{amsart} editorial wrapper at 10pt on A4, which restates the theorem-like environments the retained text needs and adds the article's own macro definitions that the retained text needs (\texttt{\textbackslash Ax}, \texttt{\textbackslash iff}, \texttt{\textbackslash then}, \texttt{\textbackslash tup}), with extra wrapper packages (bm, bbm, dsfont, xspace, cancel, mathtools, enumitem). The delivered numbering is the unit's own. Assets: no retained span contains a figure environment, a graphics-inclusion command, a PDF-page inclusion, a table or a TikZ picture; no image file is copied into this package, the package declares no asset dependency and no image file is redistributed with it. Licence: version arXiv:2608.07387v1 of this article is distributed under the Creative Commons Attribution 4.0 International licence, as recorded in the arXiv licence metadata for this exact version and shown on the abstract page https://arxiv.org/abs/2608.07387v1. A full rights notice is delivered at licenses/arxiv-2608.07387-CC-BY-4.0.txt. Characters: every retained excerpt is pure ASCII, so the delivered engine, XeLaTeX with its default Latin Modern text font, reports no missing character.
\begin{enumerate}[label=(Ax\arabic*),itemsep=1ex]
    \item $\exists x N(x)$\label{ax:zero}
    \item $\forall x (N(x) \then \exists y\ S(x,y))$\label{ax:succ}
    \item $\forall x((N(x) \wedge \neg Z(x)) \then \exists y S(y,x))$\label{ax:pred}
\end{enumerate}

\begin{lemma}\label{lem:nat}
    If $\mathcal{M},w \Vdash \text{\ref{ax:zero}--\ref{ax:pred}}$, then $\tup{N_w/\equiv,<_w^*} \cong \tup{\omega,<}$. 
    Moreover, where $\sigma\colon \tup{N_w/\equiv,<_w^*} \cong \tup{\omega,<}$, $Z_w$ is the $<_w^*$-minimal element of $N_w/\equiv$, i.e., $\sigma(Z_w) = 0$, and $S_w^* = \{\tup{a_\equiv,b_\equiv} \mid S_w(a,b)\}$ is the successor relation over $N_w/\equiv$, i.e., $S_w^*(a_\equiv,b_\equiv)$ iff $\sigma(b_\equiv) = \sigma(a_\equiv)+1$. 
\end{lemma}

\begin{enumerate}[label=(Ax\arabic*)]
\setcounter{enumi}{3}
    \item $\forall x_1x_2x_3\forall y_1y_2y_3((x_1 \equiv y_1 \wedge x_2 \equiv y_2 \wedge x_3 \equiv y_3) \then ((A(x_1,x_2,x_3) \iff A(y_1,y_2,y_3)) \wedge (M(x_1,x_2,x_3) \iff M(y_1,y_2,y_3))))$\label{ax:cong}
    % \item $\forall x\forall y((N(x) \wedge N(y)) \then (\exists z A(x,y,z) \wedge \exists z M(x,y,z)))$\label{ax:total}
    \item $\forall xyzz'(((A(x,y,z) \wedge A(x,y,z')) \then z \equiv z') \wedge ((M(x,y,z) \wedge M(x,y,z')) \then z \equiv z'))$\label{ax:func}
    \item $\forall xy((N(x) \wedge Z(y)) \then (A(x,y,x) \wedge M(x,y,y)))$\label{ax:unit}
    \item $\forall xyy'zz'((S(y,y') \wedge S(z,z')) \then (A(x,y,z) \then A(x,y',z')))$\label{ax:add}
    \item $\forall xyy'zu((S(y,y') \wedge A(z,x,u)) \then (M(x,y,z) \then M(x,y',u)))$\label{ax:prod}
\end{enumerate}

\begin{lemma}
    Let $\mathcal{M},w \Vdash \Ax$. 
    Let $\sigma\colon \tup{N_w/\equiv,<_w^*} \cong \tup{\omega,<}$. 
    Then for all $a,b,c \in N_w$:
    \begin{enumerate}
        \item $A_w(a,b,c)$ iff $\sigma(a_\equiv) + \sigma(b_\equiv) = \sigma(c_\equiv)$.
        \item $M_w(a,b,c)$ iff $\sigma(a_\equiv) \times \sigma(b_\equiv) = \sigma(c_\equiv)$.
    \end{enumerate}
\end{lemma}

\begin{proof}
    In each case, we proceed by induction on $\sigma(b_\equiv)$.
    \begin{enumerate}
        \item 
        \begin{enumerate}[label=(\roman*)]
            \item \textbf{Base Case:} Suppose $b \in N_w$ where $\sigma(b_\equiv) = 0$. 
            By Lemma~\ref{lem:nat}, $b_\equiv = Z_w$, and so $b \in Z_w$. 
            By \ref{ax:unit}, $A_w(a,b,a)$. 
            \begin{itemize}
                \item[($\Rightarrow$)] Suppose $A_w(a,b,c)$. 
                Then by \ref{ax:func}, $c \equiv_w a$, and so $c_\equiv = a_\equiv$. 
                Hence, $\sigma(a_\equiv) + \sigma(b_\equiv) = \sigma(a_\equiv) + 0 = \sigma(a_\equiv) = \sigma(c_\equiv)$. 
                
                \item[($\Leftarrow$)] Suppose $\sigma(a_\equiv) + \sigma(b_\equiv) = \sigma(c_\equiv)$. 
                Then $\sigma(a_\equiv) + 0 = \sigma(a_\equiv) = \sigma(c_\equiv)$, and so $a \equiv_w c$. 
                Since $A_w(a,b,a)$, it follows by \ref{ax:cong} that $A_w(a,b,c)$. 
            \end{itemize}

            \item \textbf{Inductive Step:} Suppose the claim holds for $b$, where $\sigma(b_\equiv) = n$. 
            Let $\sigma(b'_\equiv) = n+1$. 
            Thus, $b'_\equiv$ is the $S_w^*$-successor of $b_\equiv$, i.e., $S_w(b,b')$. 
            We'll show that the claim holds for $b'_\equiv$. 
            \begin{itemize}
                \item[($\Leftarrow$)] Suppose $\sigma(a_\equiv) + \sigma(b'_\equiv) = \sigma(c'_\equiv)$ for some $a_\equiv$ and $c'_\equiv$. 
                Since $\sigma(b'_\equiv) \neq 0$, $\sigma(c'_\equiv) \neq 0$, and so $c' \notin Z_w$ by Lemma~\ref{lem:nat}. 
                By \ref{ax:pred}, there is a $c \in N_w$ such that $S_w(c,c')$. 
                % Let $c_\equiv$ be the $S_w^*$-predecessor of $c'_\equiv$. 
                Thus, by Lemma~\ref{lem:nat}, $\sigma(a_\equiv) + \sigma(b_\equiv) = \sigma(c_\equiv)$. 
                By IH, $A_w(a,b,c)$. 
                By \ref{ax:add}, $A_w(a,b',c')$. 
                
                \item[($\Rightarrow)$] Suppose now that $\sigma(a_\equiv) + \sigma(b'_\equiv) \neq \sigma(c'_\equiv)$ for some $a_\equiv$ and $c'_\equiv$. 
                Since $\sigma$ is an isomorphism, there's a $d' \in N_w$ such that $\sigma(a_\equiv) + \sigma(b'_\equiv) = \sigma(d'_\equiv)$. 
                Thus, by the above reasoning, $A_w(a,b',d')$. 
                Since $\sigma(c'_\equiv) \neq \sigma(d'_\equiv)$, $c' \not\equiv_w d'$. 
                So by \ref{ax:func}, $\neg A_w(a,b',c')$. 
            \end{itemize}
        \end{enumerate}
     
        \item 
        \begin{enumerate}[label=(\roman*)]
            \item \textbf{Base Case:} Suppose $\sigma(b_\equiv) = 0$. 
            By Lemma~\ref{lem:nat}, $b \in Z_w$. 
            By \ref{ax:unit}, $M_w(a,b,b)$. 
            \begin{itemize}
                \item[($\Rightarrow$)] Suppose $M_w(a,b,c)$. 
                Then by \ref{ax:func}, $c \equiv_w b$, and so $\sigma(c_\equiv) = 0$. 
                Hence, $\sigma(a_\equiv) \times \sigma(b_\equiv) = \sigma(a_\equiv) \times 0 = 0 = \sigma(c_\equiv)$. 

                \item[($\Leftarrow$)] Suppose now that $\sigma(a_\equiv) \times \sigma(b_\equiv) = \sigma(c_\equiv)$. 
                Then $\sigma(c_\equiv) = \sigma(a_\equiv) \times 0 = 0$, i.e., $c \equiv_w b$. 
                Since $M_w(a,b,b)$, it follows by \ref{ax:cong} that $M_w(a,b,c)$. 
            \end{itemize}

            \item \textbf{Inductive Step:} Now suppose the claim holds for $b$ where $\sigma(b_\equiv) = n$. 
            Let $\sigma(b'_\equiv) = n+1$. 
            Thus, $b'_\equiv$ is the $S_w^*$-successor of $b_\equiv$. 
            We'll show that the claim holds for $b'_\equiv$. 
            \begin{itemize}
                \item[($\Leftarrow$)] Suppose that $\sigma(a_\equiv) \times \sigma(b'_\equiv) = \sigma(c_\equiv)$ for some $a_\equiv$ and $c_\equiv$. 
                Thus, $\sigma(c_\equiv) \geq \sigma(a_\equiv)$. 
                % If $\sigma(a_\equiv) = 0$, then $\sigma(c_\equiv) = 0 = \sigma(a_\equiv) \times \sigma(b_\equiv)$, and so $a \equiv_w c$. 
                % By \ref{ax:unit}, $M_w(a,b',a)$. 
                % So by \ref{ax:cong}, $M_w(a,b',c)$. 
                % So let $\sigma(a_\equiv) \neq 0$. 
                % Since $\sigma(b'_\equiv) \neq 0$, $\sigma(c_\equiv) > \sigma(a_\equiv)$. 
                Since $\sigma$ is an isomorphism, there is a $d$ such that $\sigma(d_\equiv) + \sigma(a_\equiv) = \sigma(c_\equiv)$. 
                (If $\sigma(a_\equiv) = \sigma(c_\equiv)$, we can let $d \in Z_w$.) 
                Thus, $\sigma(a_\equiv) \times \sigma(b_\equiv) = \sigma(d_\equiv)$. 
                By IH, $M_w(a,b,d)$. 
                By \ref{ax:prod}, $M_w(a,b',c)$. 


                \item[($\Rightarrow$)] Suppose now that $\sigma(a_\equiv) \times \sigma(b'_\equiv) \neq \sigma(c_\equiv)$ for some $a_\equiv$ and $c_\equiv$. 
                Since $\sigma$ is an isomorphism, there is a $d$ where $\sigma(a_\equiv) \times \sigma(b'_\equiv) = \sigma(d_\equiv)$. 
                Thus, by the above reasoning, $M_w(a,b',d)$. 
                Since $\sigma(c_\equiv) \neq \sigma(d_\equiv)$, $c \not\equiv_w d$. 
                So by \ref{ax:func}, $\neg M_w(a,b',c)$. 
            \end{itemize}
        \end{enumerate}
    \end{enumerate}
\end{proof}

\end{document}
